% META: {"id": "1SPE-EXPONENTIELLE-FR-R5", "chapitre": "1SPE-EXPONENTIELLE", "type_objet": "remediation", "capacites_codes": ["R5"], "status": "generated"}

%---------------------------------------------------------------
% FICHE DE REMISE A NIVEAU — R5 : Suites geometriques
% Prerequis issu de 1SPE-SUITES
%---------------------------------------------------------------

\subsection*{Rappel — Suites geometriques et coefficient multiplicateur}

Une suite $(u_n)$ est \textbf{geometrique} de raison $q$ si pour tout $n$ : $u_{n+1} = q \times u_n$.

Formule explicite : $u_n = u_0 \times q^n$.

Lien avec l'exponentielle : la suite $u_n = u_0 \times \mathrm{e}^{rn}$ est geometrique de raison $q = \mathrm{e}^r$, car $u_{n+1} = u_0 \times \mathrm{e}^{r(n+1)} = \mathrm{e}^r \times u_n$.

Le coefficient multiplicateur d'une augmentation de $t\,\%$ est $q = 1 + \dfrac{t}{100}$.

%---------------------------------------------------------------

\begin{exercice}{1SPE-EXPO-FR-R5-EX1}{1}{5}
Soit $(u_n)$ la suite geometrique de premier terme $u_0 = 5$ et de raison $q = 2$.
\begin{enumerate}
  \item Calculer $u_1$, $u_2$, $u_3$.
  \item Donner la formule explicite de $u_n$.
  \item Calculer $u_{10}$.
\end{enumerate}
\end{exercice}

\coupDePouce{1}{$u_{n+1} = q \times u_n$. Formule explicite : $u_n = u_0 \times q^n$.}

% BEGIN-VERIFY
% u0 = 5
% q = 2
% u1 = u0 * q
% assert u1 == 10
% u2 = u1 * q
% assert u2 == 20
% u3 = u2 * q
% assert u3 == 40
% assert u0 * q**10 == 5120
% END-VERIFY

\begin{corrige}{1SPE-EXPO-FR-R5-EX1}

\textbf{Question 1.}
$u_1 = 5 \times 2 = 10$, $u_2 = 10 \times 2 = 20$, $u_3 = 20 \times 2 = 40$.

\textbf{Question 2.}
$u_n = 5 \times 2^n$.

\textbf{Question 3.}
$u_{10} = 5 \times 2^{10} = 5 \times 1024 = 5120$.

\end{corrige}

%---------------------------------------------------------------

\begin{exercice}{1SPE-EXPO-FR-R5-EX2}{1}{6}
Un capital de $1\,000$~euros est place a un taux annuel de $3\,\%$. On note $C_n$ le capital au bout de $n$ annees.
\begin{enumerate}
  \item Quel est le coefficient multiplicateur annuel ?
  \item Exprimer $C_n$ en fonction de $n$.
  \item Calculer $C_5$ (arrondir au centime).
\end{enumerate}
\end{exercice}

\coupDePouce{1}{Coefficient multiplicateur pour une augmentation de $t\,\%$ : $q = 1 + t/100$. Ici $t = 3$.}

% BEGIN-VERIFY
% from sympy import Rational, N
% q = Rational(103, 100)
% C0 = 1000
% C5 = C0 * q**5
% assert abs(float(C5) - 1159.27) < 0.01
% END-VERIFY

\begin{corrige}{1SPE-EXPO-FR-R5-EX2}

\textbf{Question 1.} Le coefficient multiplicateur est $q = 1 + \frac{3}{100} = 1{,}03$.

\textbf{Question 2.} $(C_n)$ est geometrique de raison $1{,}03$ et $C_0 = 1000$.
\[
C_n = 1000 \times 1{,}03^n.
\]

\textbf{Question 3.}
$C_5 = 1000 \times 1{,}03^5 \approx 1000 \times 1{,}15927 \approx 1159{,}27$~euros.

\end{corrige}

%---------------------------------------------------------------

\begin{exercice}{1SPE-EXPO-FR-R5-EX3}{2}{6}
Une population de bacteries est modelisee par $P_n = 500 \times (1{,}5)^n$ ou $n$ est le nombre d'heures.
\begin{enumerate}
  \item Justifier que $(P_n)$ est une suite geometrique. Donner sa raison.
  \item Calculer $P_0$, $P_1$, $P_2$ et $P_3$.
  \item A partir de quel rang $n$ a-t-on $P_n > 5\,000$ ? (on pourra tester quelques valeurs)
\end{enumerate}
\end{exercice}

\coupDePouce{1}{$P_{n+1}/P_n = (1{,}5)^{n+1}/(1{,}5)^n = 1{,}5$ : le quotient est constant.}

% BEGIN-VERIFY
% from sympy import Rational
% q = Rational(3, 2)
% P = [500 * q**n for n in range(10)]
% assert P[0] == 500
% assert P[1] == 750
% assert P[2] == Rational(2250, 2)  # 1125
% assert float(P[2]) == 1125.0
% assert float(P[3]) == 1687.5
% # P_n > 5000 : tester
% n = 0
% while float(500 * q**n) <= 5000:
%     n += 1
% assert n == 6  # P_6 = 500 * 1.5^6 = 500 * 11.390625 = 5695.3125
% END-VERIFY

\begin{corrige}{1SPE-EXPO-FR-R5-EX3}

\textbf{Question 1.} $P_{n+1} = 500 \times (1{,}5)^{n+1} = 1{,}5 \times 500 \times (1{,}5)^n = 1{,}5 \times P_n$.

Le quotient $P_{n+1}/P_n = 1{,}5$ est constant : $(P_n)$ est geometrique de raison $q = 1{,}5$.

\textbf{Question 2.}
$P_0 = 500$, $P_1 = 750$, $P_2 = 1125$, $P_3 = 1687{,}5$.

\commentaireMarge{On teste les valeurs successives jusqu'a depasser $5000$.}

\textbf{Question 3.} On calcule :
$P_4 = 2531{,}25$, $P_5 \approx 3796{,}9$, $P_6 \approx 5695{,}3$.

$P_5 \approx 3797 < 5000$ et $P_6 \approx 5695 > 5000$.

A partir du rang $n = 6$, on a $P_n > 5000$.

\end{corrige}
